\chapter{Grandeurs, aires et volumes}
\label{exercices:17}
\section*{Exercice 17.1 -- Conversions}
Convertir $0{,}45\unite{m^2}$ en $\unite{cm^2}$, $3250\unite{cm^3}$ en litres, $2\unite h\,18\unite{min}$ en heures décimales et $90\unite{km/h}$ en $\unite{m/s}$.

\section*{Exercice 17.2 -- Aires composées}
Calculer l'aire d'un trapèze de bases $8$ et $14\unite{cm}$, de hauteur $5\unite{cm}$. Calculer ensuite l'aire exacte d'une couronne de rayons $7$ et $4\unite{cm}$. Justifier les opérations.

\section*{Exercice 17.3 -- Assemblage et débit}
Un objet plein est formé d'un cylindre de rayon $3\unite{cm}$ et de hauteur $8\unite{cm}$, surmonté d'une demi-boule de même rayon. Calculer son volume exact. Un réservoir de $1200\unite L$ se remplit à $25\unite{L/min}$ : combien de temps pour atteindre $75\,\%$ de sa capacité, en partant de zéro ?

\section*{Exercice 17.4 -- Grandeurs composées}
On parcourt $30\unite{km}$ à $60\unite{km/h}$ puis $30\unite{km}$ à $30\unite{km/h}$. Calculer la vitesse moyenne. Un liquide de masse volumique $0{,}8\unite{kg/L}$ occupe $5\unite L$ : calculer sa masse.
